% Target: rmp-ii-peps-rg
% Formula: A renormalized 2x2 PEPS plaquette decomposes into the four source panels
%   (a)-(d).
% Ink: Canonical public tenkzlattice, tenkz environments; no raw TikZ.
% Boundary: Each PEPS factor keeps its exterior legs open; the final bead closes only
%   its virtual ring.
% Source: paper TN-Review-main.tex line 1025; author source ImagesReview Section
%   II/ImagesReview Section II.tex lines 742-800
% Capabilities: grid, lattice, typed-ports
(a)
\[
  \begin{tenkzlattice}[rows=2, cols=2, frame=oblique, boundary=open,
                       physical=up]
    \tnsite{(1,1)}
\end{tenkzlattice}
\]

(b)
\[
  \begin{tenkzlattice}[rows=1, cols=1, frame=oblique, boundary=open,
                       physical=up]
    \tnsite{(1,1)}
\end{tenkzlattice}
  \;\otimes\;
  % three beads with depth (a sheet): up+south / west+up+north / up+east,
  % matching the source's three oblique corner tensors
  \begin{tenkzlattice}[rows=1, cols=1, frame=oblique, boundary=open,
                       physical=up, west=none, east=none, north=none]
    \tnsite{(1,1)}
\end{tenkzlattice}
  \;
  \begin{tenkzlattice}[rows=1, cols=1, frame=oblique, boundary=open,
                       physical=up, east=none, south=none]
    \tnsite{(1,1)}
\end{tenkzlattice}
  \;
  \begin{tenkzlattice}[rows=1, cols=1, frame=oblique, boundary=open,
                       physical=up, west=none, north=none, south=none]
    \tnsite{(1,1)}
\end{tenkzlattice}
\]

(c)
\[
  \begin{tenkzlattice}[rows=1, cols=1, frame=oblique, boundary=open,
                       physical=up]
    \tnsite{(1,1)}
\end{tenkzlattice}
  \;\otimes\;
  % three independent beads with depth (a sheet): west+up / up+south /
  % up+north, matching the source's three separate oblique corner tensors
  \begin{tenkzlattice}[rows=1, cols=1, frame=oblique, boundary=open,
                       physical=up, east=none, north=none, south=none]
    \tnsite{(1,1)}
\end{tenkzlattice}
  \;
  \begin{tenkzlattice}[rows=1, cols=1, frame=oblique, boundary=open,
                       physical=up, west=none, east=none, north=none]
    \tnsite{(1,1)}
\end{tenkzlattice}
  \;
  \begin{tenkzlattice}[rows=1, cols=1, frame=oblique, boundary=open,
                       physical=up, west=none, east=none, south=none]
    \tnsite{(1,1)}
\end{tenkzlattice}
\]

(d)
\[
  \begin{tenkzlattice}[rows=1, cols=1, frame=oblique, boundary=open,
                       physical=up]
    \tnsite{(1,1)}
\end{tenkzlattice}
  \;\otimes\;
  \begin{tenkzlattice}[rows=1, cols=1, frame=oblique, boundary=open,
                       physical=up]
    \tnsite{(1,1)}
\end{tenkzlattice}
  \;\otimes\;
  % the bead whose virtual wire closes on itself: hooks trace closure,
  % opening upward as the source draws it (frame rotated 180 degrees, with
  % physical=down pre-compensating so the physical leg still points up)
  \begin{tenkz}[physical=down, boundary=periodic, trace style=hooks,
                frame={rotate=180}]
    \tn{}
  \end{tenkz}
\]
